\begin{theorem}[Wild odd boundary pencil]\label{mod:cases-wildodd-boundarypencil}
Let $F$ be a nonarchimedean local field, let $p\ne2$ be prime with
$\operatorname{char}(k_F)=p$, and let $C$ be residual additive-character
data whose lower character $\psi_0$ is nontrivial and Frobenius invariant.
Let
\[
 (\eta_0,\eta_0\gamma_0),\qquad (\eta,\eta\gamma)
\]
be respectively the certified boundary generator and base rows from
\texttt{ResidualCoefficients}, and put
\[
 d_0=\eta/\eta_0,\qquad
 \rho=d_0^{-1},\qquad
 \sigma=\gamma,\qquad
 \tau=\rho\gamma_0.
\]
Assume the imported nonzero-index boundary facts
$[j]+d_0\ne0$ for every $j\in\F_p^\times$.  With the exact upstream
inverse-Frobenius choices
\[
 a_0=\operatorname{Frob}^{-1}(1-d_0^{p-1}),
 \qquad
 b_0=\operatorname{Frob}^{-1}(\gamma-\gamma_0),
\]
one has the boundary upper--norm versus twist identity
\[
 Q_{\psi_0}(a_0,b_0)
 \prod_{u\in\F_p^\times}Q_{\psi_0}([u]\rho,[u]\tau)
 =
 \prod_{j\in\F_p}
 Q_{\psi_0}(1+[j]\rho,\sigma+[j]\tau).
\]
Here the prime-field scalars are exactly the values of
\texttt{ZMod.castHom}; the left product is indexed by
$(\operatorname{ZMod}p)^\times$, while the right product includes the zero
index.  Thus the upper factor is on the left and every boundary twist factor
is on the right, with the Frobenius direction and both affine coefficients
unchanged from the residual-coefficient certificate.
\LeanFile{LanglandsFirstMainLemma/Cases/WildOdd/BoundaryPencil.lean}
\lean{LanglandsFirstMainLemma.wildOdd_boundaryPencil}
\leanok
\SourceText{Lemma lem:affine-pencil and the boundary part of Proposition prop:odd-phase-cancel.}
\uses{mod:finitefield-affinepencil,mod:cases-wildodd-residualcoefficients}
\FormalizationContract
The proof constructs
\texttt{wildOdd\_boundaryAffinePencilCertificate} from the two named rows
and the nonzero-index facts, then applies \texttt{affinePencil} directly to
its four fields.  It performs no residual-coefficient calculation, no
coordinate transport, no scalar-phase product, and no parity reduction.
\end{theorem}
\begin{proof}
The boundary certificate supplies $\rho\ne0$, all conditions
$1+[j]\rho\ne0$, and the two exact identities
\[
 a_0^p=1-(\rho^{p-1})^{-1},
 \qquad
 b_0^p=\sigma-\rho^{-1}\tau.
\]
Apply Lemma~\ref{lem:affine-pencil} with these data and with the Frobenius
invariance carried by $C$.
\end{proof}
